\apendice{Dicionário de variáveis}
\label{ap:variaveis}

\begin{orientacao}[Dicionário de variáveis]
  O dicionário de variáveis documenta cada atributo da base utilizada na modelagem e permite que outro pesquisador compreenda, audite e reproduza o estudo. Liste todas as variáveis que entram no modelo, inclusive as derivadas (defasagens, médias móveis, indicadores calculados), e identifique a variável-alvo. Para cada uma, informe o tipo, a unidade ou o domínio de valores, a origem e, no caso das variáveis derivadas, a regra de cálculo. Essa informação é decisiva para afastar o vazamento de dados: uma variável derivada que use informação posterior ao instante da previsão invalida a avaliação do modelo. Se a base tiver muitas variáveis, agrupe-as por dimensão (temporal, meteorológica, operacional) ou divida o quadro em mais de uma página. No texto, cite o quadro antes de ele aparecer.
\end{orientacao}

O Quadro~\ref{qua:variaveis} descreve as variáveis da base de dados utilizada na modelagem, com a origem de cada uma. As variáveis derivadas foram calculadas apenas com informações disponíveis até o instante da previsão.

\begin{quadros}[H]
  \centering
  \caption{Dicionário de variáveis da base horária}
  \label{qua:variaveis}
  \small
  \singlespacing
  \setlength{\tabcolsep}{4pt}
  \renewcommand{\arraystretch}{1.3}
  \begin{tabular}{|l|c|c|l|l|}
    \hline
    \textbf{Variável} & \textbf{Tipo} & \textbf{Unidade} & \textbf{Descrição} & \textbf{Origem} \\ \hline
    \texttt{lentidao\_km}   & contínua   & km       & Extensão total de lentidão (alvo)    & CET-SP \\ \hline
    \texttt{lentidao\_lag1} & contínua   & km       & Lentidão observada na hora anterior  & Derivada \\ \hline
    \texttt{hora}           & discreta   & 0 a 23   & Hora do dia                          & Derivada \\ \hline
    \texttt{dia\_semana}    & categórica & 1 a 7    & Dia da semana                        & Derivada \\ \hline
    \texttt{feriado}        & binária    & 0 ou 1   & Indicador de feriado                 & Calendário oficial \\ \hline
    \texttt{precipitacao}   & contínua   & mm       & Precipitação acumulada na hora       & INMET \\ \hline
    \texttt{temperatura}    & contínua   & °C       & Temperatura média do ar              & INMET \\ \hline
    \texttt{ocorrencias}    & discreta   & contagem & Ocorrências viárias ativas           & CET-SP \\ \hline
    \texttt{eventos}        & binária    & 0 ou 1   & Evento de grande porte na cidade     & Agenda municipal \\ \hline
  \end{tabular}

  \fonte{Elaborado pela autora (\ano).}
  {\footnotesize\singlespacing Nota: CET-SP: Companhia de Engenharia de Tráfego de São Paulo; INMET: Instituto Nacional de Meteorologia. Variáveis derivadas: \texttt{lentidao\_lag1} corresponde a \texttt{lentidao\_km} defasada em uma hora; \texttt{hora} e \texttt{dia\_semana} foram extraídas do carimbo de data e hora do registro.\par}
\end{quadros}